% Part: lambda-calculus
% Chapter: syntax
% Section: de-bruijn

\documentclass[../../../include/open-logic-section]{subfiles}

\begin{document}

\olfileid[fr]{lam}{syn}{deb}

\olsection{Indices de de Bruijn}

L'$\alpha$-équivalence est très naturelle :
des termes $\alpha$-équivalents ont le
« même sens ». On peut même définir une
syntaxe des lambda-termes qui ne distingue
pas les termes que l'on peut convertir
l'un en l'autre par $\alpha$-conversion.
La représentation la plus connue remplace
les variables par leurs \emph{indices de
de Bruijn}.

Dans $\lambd[x][M]$, la variable~$x$
nomme explicitement le paramètre de la
fonction; on peut donc l'employer dans~$M$
pour y faire référence. Avec les indices de
de Bruijn, les paramètres n'ont plus de
nom : une occurrence dans le corps reçoit
un nombre qui indique combien d'abstractions
la séparent de celle qui la lie. Considérons
par exemple $\lambd[x][\lambd[y][y x]]$.
L'abstraction extérieure lie~$x$ et
l'abstraction intérieure lie~$y$. Le
sous-terme $y x$ est dans la portée de
l'abstraction intérieure : aucune abstraction
ne sépare $y$ de son $\lambd[y]$, tandis
qu'une abstraction sépare $x$ de son
$\lambd[x]$. On écrit donc $0\, 1$
pour $y x$, et $\lambd[][\lambd[][0\, 1]]$
pour le terme entier.

\begin{explain}
L'espace dans $0\,1$ sépare les deux indices d'une application.
Chaque indice peut lui-même s'écrire avec plusieurs chiffres.
\end{explain}

\begin{defn}
Les termes de de Bruijn se définissent
inductivement comme suit :
\begin{enumerate}
\item $n$, où $n$ est un entier naturel.
\item $PQ$, où $P$ et $Q$ sont des termes
  de de Bruijn.
\item $\lambd[][N]$, où $N$ est un terme
  de de Bruijn.
\end{enumerate}
\end{defn}

La traduction formelle des lambda-termes
ordinaires en termes indicés est la suivante.
\begin{defn}
\begin{align*}
  F_\Gamma(x) &= \Gamma(x) \\
  F_\Gamma(PQ) &= F_\Gamma(P)F_\Gamma(Q) \\
  F_\Gamma(\lambd[x][N])
    &= \lambd[][F_{x,\Gamma}(N)]
\end{align*}
Ici, $\Gamma$ est une liste de variables
indicée à partir de zéro, et $\Gamma(x)$
désigne la position de la première occurrence
de $x$ dans $\Gamma$. Par exemple, si
$\Gamma$ est la liste $x,y,z$,
$\Gamma(x)$ vaut $0$ et $\Gamma(z)$ vaut $2$.

La notation $x,\Gamma$ désigne la liste
obtenue en plaçant $x$ en tête de $\Gamma$.
Ainsi, si $\Gamma$ est la liste précédente,
$w,\Gamma$ est la liste $w,x,y,z$.
La traduction est définie lorsque toute
variable libre du terme figure dans la liste
initiale $\Gamma$.
\end{defn}

\begin{explain}
Le contexte se lit depuis le lieur le plus proche. Par exemple,
dans $\lambd[x][\lambd[y][\lambd[x][x y]]]$, le dernier
$x$ reçoit l'indice~$0$ et $y$ reçoit l'indice~$1$.
Le contexte du corps est la liste $x,y,x$.
\end{explain}

On retrouve un lambda-terme ordinaire à
partir d'un terme de de Bruijn comme suit.

\begin{defn}
\begin{align*}
  G_\Gamma(n) &= \Gamma[n] \\
  G_\Gamma(PQ) &= G_\Gamma(P) G_\Gamma(Q) \\
  G_\Gamma(\lambd[][N])
    &= \lambd[x][G_{x,\Gamma}(N)]
\end{align*}
Ici encore, $\Gamma$ est une liste de
variables indicée à partir de zéro, et
$\Gamma[n]$ désigne la variable en
position~$n$. Par exemple, si $\Gamma$
est la liste $x,y,z$, alors $\Gamma[1]$
est $y$.

Dans la dernière équation, on choisit
une variable $x$ absente de $\Gamma$.
\end{defn}

\begin{explain}
Le décodage n'est défini que si chaque indice
rencontré possède une entrée dans le contexte
courant.
Une abstraction prolonge ce contexte d'une entrée. Pour
retrouver exactement les indices d'un terme en le décodant
puis en le recodant, il suffit que le contexte initial
contienne des noms distincts. Avec le contexte $x,x$, par
exemple, l'indice~$1$ se décode en $x$, qui se recode en~$0$.
\end{explain}

Nous énonçons quelques résultats sans les
démontrer.

\begin{prop}
Si $M \aconvone M'$ et si $\Gamma$ est une
liste contenant $\FV{M}$, alors
$F_\Gamma(M) \eqs F_\Gamma(M')$.
\end{prop}

\end{document}
